\section{The extended term}
\label{sec:6.7}
Everything so far runs on the Charter as courts read it, plus the corroborated licence. The extended term adds three more departures, and the rest of the book needs only what each says. Appendix~\ref{sec:F} gives the argument for each and works through the cases.

\emph{Standing to defend.} Article 51 speaks only of an armed attack against a Member of the United Nations. The floor extends the emergency licence to any authority that administers territory under the rule in \S\ref{sec:6.2}, Member or not, and to the forces of a population a disinterested body has found occupied. Whether a population may be defended shouldn't turn on a seat other states vote on.

\emph{The genocide exception.} The term lifts for force to stop a genocide under way, even where the Council is blocked, on a merits judgment of the International Court of Justice under the Genocide Convention, or on two findings of different kinds, at least one from a body that took evidence. For this exception alone, a provisional-measures order finding the claim plausible and the risk real counts as one of the two; a Council or Assembly resolution corroborates, and two votes never make the pair. The exception lifts only for the force and place the finding covers: against the forces committing the genocide, where they are committing it, the least that stops it. Force elsewhere against the state responsible, on the theory that the cost will make it stop, is outside it. A finding under this exception also lifts the named belligerent's own emergency licence for the force and place it covers (Appendix~\ref{sec:F.4}).

\emph{The occupation exception.} A finding that territory is occupied gives its population's forces the defensive licence above, and no more. A finding that the occupation is unlawful and must end licenses force to end it: against the occupier's forces, inside the occupied territory, the least that ends it there, and nothing that changes who lives there. The evidence is the genocide exception's: an opinion or judgment of the Court, or two findings of different kinds, one of them from a body that took evidence. Cyprus and Nagorno-Karabakh are the cases (Appendix~\ref{sec:F.5}).

Standing to defend and the two exceptions are departures from the instruments, argued to the bar every departure must clear.
